\section{Second Pass: Definitions Under a Microscope}
\chapterline

\subsection{From physical intuition to formal axiomatics: Euclid to Hilbert}

Chapter~28 showed how a quantifier swap can break a proof. But logical syntax is only part of the foundation. A statement like $\forall x \in A, \; P(x)$ depends on how $A$ and $P(x)$ are defined. If definitions rest on informal mental pictures rather than explicit rules, unstated assumptions slip in unnoticed.

Euclid's \emph{Elements} (c.~300~BCE) introduced the axiomatic method to mathematics, yet Euclid opened his work with intuitive descriptions rather than formal tests: ``A point is that which has no part''; ``A line is breadthless length.''

These were appeals to visual imagination rather than operational definitions. Because Euclid relied on physical intuition for what points and lines ``were,'' unstated assumptions entered his proofs. In his first proposition, constructing an equilateral triangle, Euclid drew two overlapping circles and assumed without an axiom that they intersect. In later proofs, he assumed that a line entering a triangle through one edge must exit through another (a fact that requires Pasch's axiom).

For centuries, mathematicians treated these gaps as minor because everyone agreed on the geometry of physical space. But in the nineteenth century, the discovery of non-Euclidean geometries showed that alternative, consistent geometric worlds exist. On a sphere, every pair of lines crosses; on a hyperbolic surface, infinitely many parallel lines pass through a single point. Which geometry was physically true?

In 1899, David Hilbert published \emph{Grundlagen der Geometrie}, redefining mathematical definitions: not descriptions of physical substance, but formal relations that abstract objects satisfy.

\subsection{Definitions as precise classification tests}

In everyday conversation, words have soft edges. If someone asks you to define a ``chair'' or a ``game,'' you might describe typical examples: something with four legs you sit on, or an activity with rules and a winner. But if you encounter a beanbag, an office stool, or solitaire, the colloquial boundary blurs. Human language tolerates ambiguity because context fills the gaps.

Pure mathematics cannot rely on fuzzy boundaries. A mathematical definition operates by an explicit set of conditions rather than typical examples or sensory impressions. It specifies an exact checklist of requirements: an object either satisfies every condition and qualifies, or fails at least one and does not. There is no intermediate case.

\begin{center}
\begin{tikzpicture}[
  >=Latex,
  node distance=0.8cm,
  comp/.style={draw=#1!80!black, fill=#1!12, thick, rounded corners=3pt, font=\small, text width=6.6cm, align=left, inner sep=6pt},
  flow/.style={->, very thick, draw=gray!70}
]

  % Title
  \node[font=\bfseries, text=structurepurple] at (0, 1.2) {The Four Anatomical Layers of a Definition};

  % Component 1: Ambient Context
  \node[comp=sourceblue] (c1) at (0, 0) {
    \textbf{1. Ambient Universe (Prerequisites)}\\
    Sets the stage and valid operations.\\
    \emph{Example}: Let $f: A \to B$ be a function between sets.
  };

  % Component 2: Input Signature
  \node[comp=quotientorange, below=0.7cm of c1] (c2) {
    \textbf{2. Input Signature (Subject)}\\
    Identifies the specific candidate object being tested.\\
    \emph{Example}: The function $f$ itself.
  };

  % Component 3: Defining Condition
  \node[comp=structurepurple, below=0.7cm of c2] (c3) {
    \textbf{3. Defining Contract (The Test)}\\
    The exact logical predicate that must be verified.\\
    \emph{Example}: $\forall a_1, a_2 \in A, \; f(a_1) = f(a_2) \implies a_1 = a_2$.
  };

  % Component 4: Downstream Rights
  \node[comp=targetgreen, below=0.7cm of c3] (c4) {
    \textbf{4. Downstream Rights (Consequences)}\\
    Guaranteed properties and algebraic tools gained.\\
    \emph{Example}: Left inverse exists; information is never destroyed.
  };

  % Arrows connecting components
  \draw[flow] (c1) -- (c2);
  \draw[flow] (c2) -- (c3);
  \draw[flow] (c3) -- (c4);

\end{tikzpicture}\\[0.4em]
{\small\textbf{Figure 29.1:} The four anatomical layers of a mathematical definition. A rigorous definition establishes an ambient universe, identifies the candidate input, tests it against a logical defining condition, and provides concrete consequences once the condition is verified.}
\end{center}

\subsection{The four anatomical parts of a definition}

Whenever you meet a new mathematical definition, resist the temptation to memorize the sentence as a long string of words. Dissect it into its four architectural layers:

\begin{enumerate}
\item \textbf{The Ambient Universe}: What mathematical world are we standing in? Does this definition require a simple set of numbers, an algebraic group (a collection of symmetries with a reversible operation), a metric space (a set equipped with a notion of distance between points), or a vector space (a space where objects can be stretched and added together)? If the ambient prerequisites are missing, the definition is completely untethered.
\item \textbf{The Input Signature}: What object is receiving the descriptive label? Is the definition classifying a single point, a subset, an ordered pair, or a function?
\item \textbf{The Defining Contract}: What exact condition must be checked? Is it an algebraic equation, an if-then implication, or an adversarial game with quantifiers?
\item \textbf{The Downstream Rights}: What new deductive tools do you gain once an object passes the test? If a set is open in a metric space, you can find an open ball around any of its points that remains entirely within the set. If a function is injective (one-to-one), you gain the right to cancel $f$ from both sides of an equation: $f(a_1) = f(a_2)$ implies $a_1 = a_2$.
\end{enumerate}

\begin{recoverybox}
\textbf{Pause and think}: Suppose someone asks you: ``Is the square root function injective?'' Before attempting a proof, check Layer~1 (the ambient universe). If the function is defined on the non-negative real numbers, $f: [0, \infty) \to \R$, it is strictly increasing and passes the test. But if someone attempts to define $f$ on all real numbers without restricting the domain, the input $\sqrt{-4}$ is not even a real number! Without establishing the ambient universe first, testing the defining contract produces an immediate type error.
\end{recoverybox}

\subsection{Microscope study: why injective is defined via equality}

Consider the definition of an injective (one-to-one) function. The informal intuition is simple:
\begin{quote}
``Different inputs produce different outputs.''
\end{quote}
Translated directly into logic, this intuition reads:
\[
\forall a_1, a_2 \in A, \; a_1 \neq a_2 \implies f(a_1) \neq f(a_2).
\]
Yet almost every textbook in higher mathematics chooses instead to define injectivity using its \textbf{contrapositive}---the logically equivalent statement formed by reversing the implication and negating both sides ($P \implies Q$ becomes $\text{not } Q \implies \text{not } P$):
\[
\forall a_1, a_2 \in A, \; f(a_1) = f(a_2) \implies a_1 = a_2.
\]

Why prefer the contrapositive over the intuitive sentence?

The reason comes down to practical proof mechanics: \textbf{equations are vastly easier to compute with than inequalities.}
\begin{itemize}
\item If you start with $a_1 \neq a_2$, what can you do algebraically? Very little. Inequalities do not play well with standard algebra. If you square two unequal numbers, they can easily become equal ($(-3)^2 = 3^2 = 9$). If you multiply both sides by an unknown variable $z$, both sides collapse to zero if $z = 0$. An inequality merely tells you that two things are not identical; it offers no common handle to grab onto.
\item If you start with $f(a_1) = f(a_2)$, you have an equation. Equality allows direct substitution: wherever $f(a_1)$ appears, it can be replaced by $f(a_2)$. You can apply algebraic identities, factor both sides, cancel non-zero terms, and systematically isolate the variables until you arrive at $a_1 = a_2$.
\end{itemize}

\begin{example}
Prove that $f: \R \to \R$ given by $f(x) = 5x - 7$ is injective.
\begin{proof}
Let $a_1, a_2 \in \R$ be arbitrary, and assume $f(a_1) = f(a_2)$. By definition of $f$,
\[
5a_1 - 7 = 5a_2 - 7.
\]
Adding $7$ to both sides yields $5a_1 = 5a_2$. Dividing both sides by $5$ gives $a_1 = a_2$. Therefore, $f$ is injective.
\end{proof}
Notice how direct and effortless the algebraic chain is when built on equalities.
\end{example}

\begin{recoverybox}
\textbf{Pause and think}: A beginner often tries to verify that a function is injective by assuming $a_1 = a_2$ and showing that $f(a_1) = f(a_2)$. Why is this completely wrong? The statement $a_1 = a_2 \implies f(a_1) = f(a_2)$ is not injectivity at all---it is the definition of a \emph{well-defined function} (equal inputs give equal outputs). Every function satisfies that property, even the constant function $f(x) = 42$! Injectivity requires running the arrow in reverse: equal outputs must force equal inputs.
\end{recoverybox}

\begin{center}
\begin{tikzpicture}[
  >=Latex,
  card/.style={draw=gray!50, fill=gray!5, rounded corners=5pt, thick},
  valnode/.style={circle, fill=sourceblue, inner sep=1.6pt},
  badval/.style={circle, fill=trapred, inner sep=1.6pt},
  goodval/.style={circle, fill=targetgreen!80!black, inner sep=1.6pt}
]
  % Title
  \node[font=\bfseries, text=structurepurple] at (0, 3.2) {Well-Defined Function versus Injective Function};

  % --- LEFT PANEL: WELL-DEFINED ---
  \node[font=\small\bfseries, text=sourceblue!90!black] at (-3.7, 2.6) {(a) Well-Defined: $a_1 = a_2 \implies f(a_1) = f(a_2)$};
  \node[font=\scriptsize, text=gray!80!black] at (-3.7, 2.2) {Forward Determinism (No Output Splitting)};

  % Box for panel A
  \draw[draw=sourceblue!40, fill=sourceblue!4, rounded corners=6pt] (-6.8, -2.4) rectangle (-0.6, 1.9);

  % Valid function: 1 input -> 1 output
  \node[goodval, label={left:\scriptsize $a$}] (wa1) at (-5.6, 1.0) {};
  \node[goodval, label={right:\scriptsize $f(a)$}] (wb1) at (-1.8, 1.0) {};
  \draw[->, thick, targetgreen!80!black] (wa1) -- (wb1);
  \node[font=\tiny\bfseries, text=targetgreen!80!black] at (-3.7, 1.3) {Single output \checkmark};

  % Forbidden: 1 input -> 2 outputs
  \node[badval, label={left:\scriptsize $a$}] (wa2) at (-5.6, -0.6) {};
  \node[badval, label={right:\scriptsize $b_1$}] (wb2a) at (-1.8, -0.1) {};
  \node[badval, label={right:\scriptsize $b_2$}] (wb2b) at (-1.8, -1.1) {};
  \draw[->, thick, dashed, trapred] (wa2) -- (wb2a);
  \draw[->, thick, dashed, trapred] (wa2) -- (wb2b);
  \node[font=\tiny\bfseries, text=trapred] at (-3.5, -0.6) {Forbidden: Splitting \texttimes};

  \node[font=\scriptsize, text=gray!80!black, align=center] at (-3.7, -1.9) {
    \textbf{All functions} must be well-defined:\\
    identical inputs must give identical outputs.
  };

  % --- RIGHT PANEL: INJECTIVE ---
  \node[font=\small\bfseries, text=structurepurple!90!black] at (3.7, 2.6) {(b) Injective: $f(a_1) = f(a_2) \implies a_1 = a_2$};
  \node[font=\scriptsize, text=gray!80!black] at (3.7, 2.2) {Backward Reflection (No Collisions)};

  % Box for panel B
  \draw[draw=structurepurple!40, fill=structurepurple!4, rounded corners=6pt] (0.6, -2.4) rectangle (6.8, 1.9);

  % Injective: 2 distinct inputs -> 2 distinct outputs
  \node[goodval, label={left:\scriptsize $a_1$}] (ia1) at (1.8, 1.3) {};
  \node[goodval, label={left:\scriptsize $a_2$}] (ia2) at (1.8, 0.7) {};
  \node[goodval, label={right:\scriptsize $f(a_1)$}] (ib1) at (5.6, 1.3) {};
  \node[goodval, label={right:\scriptsize $f(a_2)$}] (ib2) at (5.6, 0.7) {};
  \draw[->, thick, targetgreen!80!black] (ia1) -- (ib1);
  \draw[->, thick, targetgreen!80!black] (ia2) -- (ib2);
  \node[font=\tiny\bfseries, text=targetgreen!80!black] at (3.7, 1.6) {Distinct outputs \checkmark};

  % Forbidden for injection: 2 distinct inputs -> 1 collapsed output
  \node[badval, label={left:\scriptsize $a_1$}] (ja1) at (1.8, -0.1) {};
  \node[badval, label={left:\scriptsize $a_2$}] (ja2) at (1.8, -1.1) {};
  \node[badval, label={right:\scriptsize $f(a_1) = f(a_2)$}] (jb) at (5.2, -0.6) {};
  \draw[->, thick, dashed, trapred] (ja1) -- (jb);
  \draw[->, thick, dashed, trapred] (ja2) -- (jb);
  \node[font=\tiny\bfseries, text=trapred] at (3.3, -0.6) {Collision: Non-Injective \texttimes};

  \node[font=\scriptsize, text=gray!80!black, align=center] at (3.7, -1.9) {
    \textbf{Only injective functions} avoid collisions:\\
    equal outputs force equal inputs.
  };
\end{tikzpicture}\\[0.4em]
{\small\textbf{Figure 29.2:} Forward determinism versus backward reflection. (a) A relation is a well-defined function if every input produces a unique output ($a_1 = a_2 \implies f(a_1) = f(a_2)$); output splitting is forbidden. (b) A function is injective if distinct inputs never collide ($f(a_1) = f(a_2) \implies a_1 = a_2$); output collapsing is forbidden. Beginners often prove forward well-definedness when asked to prove injectivity.}
\end{center}

\subsection{Stress-testing definitions: dropping an axiom}

The clearest way to understand why a definition contains the exact clauses it does is to ask: \emph{what breaks if one condition is omitted?}

Definitions in pure mathematics rarely include redundant conditions. Each clause serves a specific deductive purpose. Omitting an axiom broadens the class of objects under study, often admitting structures that lack the key properties of the original concept.

\begin{example}
Consider the definition of an \textbf{equivalence relation} on a set $S$---a relation that acts like generalized equality, partitioning elements into non-overlapping equivalence classes:
\begin{enumerate}
\item \textbf{Reflexivity}: $\forall x \in S, \; x \sim x$ (every element is related to itself).
\item \textbf{Symmetry}: $\forall x, y \in S, \; x \sim y \implies y \sim x$ (the relation is symmetric between pairs).
\item \textbf{Transitivity}: $\forall x, y, z \in S, \; (x \sim y \land y \sim z) \implies x \sim z$ (the relationship extends across intermediate elements).
\end{enumerate}

What happens if we drop \textbf{transitivity}, keeping only reflexivity and symmetry? A relation satisfying only reflexivity and symmetry is called a \emph{tolerance relation}.

Consider the relation on the real number line given by $|x - y| \le 1$ (``is within $1$ unit of''):
\begin{itemize}
\item Every number is within $1$ unit of itself: $|x - x| = 0 \le 1$ (reflexive).
\item Distance is symmetric: $|x - y| = |y - x| \le 1$ (symmetric).
\item But consider $x = 0$, $y = 0.7$, and $z = 1.4$. Here $|0 - 0.7| = 0.7 \le 1$ (so $0 \sim 0.7$) and $|0.7 - 1.4| = 0.7 \le 1$ (so $0.7 \sim 1.4$). Yet $|0 - 1.4| = 1.4 > 1$, so $0 \not\sim 1.4$. Transitivity fails.
\end{itemize}

Why is transitivity so vital? Because without transitivity, a set cannot be \textbf{partitioned}---sliced cleanly into mutually disjoint subsets that cover the entire universe without overlapping. 

In a genuine equivalence relation, every element belongs to an \textbf{equivalence class} $[x] = \{y \in S : x \sim y\}$. Transitivity guarantees an all-or-nothing property: for any two elements $a$ and $b$, either their classes are completely identical ($[a] = [b]$) or they do not share any elements ($[a] \cap [b] = \emptyset$). This allows mathematicians to construct a \textbf{quotient set} $S/{\sim}$, viewing each entire class as a single object (the way all equivalent fractions $\frac{1}{2}, \frac{2}{4}, \frac{3}{6}$ are identified with the single rational number $0.5$).

Without transitivity, the neighborhood around $0$ contains $0.7$, and the neighborhood around $1.4$ also contains $0.7$, but the two neighborhoods do not coincide. Without transitivity, the equivalence classes fail to partition the set, and one cannot form a quotient space.
\end{example}

\begin{center}
\begin{tikzpicture}[
  >=Latex,
  partition/.style={draw=sourceblue!70!black, fill=sourceblue!8, rounded corners=6pt, thick},
  point/.style={circle, fill=sourceblue!90!black, inner sep=1.8pt},
  badpoint/.style={circle, fill=trapred!90!black, inner sep=1.8pt},
  rel/.style={<->, thick, sourceblue!80!black},
  failrel/.style={<->, thick, dashed, trapred!80!black}
]

  % Title
  \node[font=\bfseries, text=structurepurple] at (0, 3.2) {Equivalence Partitions versus Overlapping Tolerance};

  % Left Panel: With Transitivity
  \node[font=\bfseries\small, text=sourceblue!90!black] at (-3.4, 2.6) {With Transitivity: Disjoint Classes};

  % Class 1
  \draw[partition] (-5.3, -0.4) rectangle (-3.5, 2.1);
  \node[font=\scriptsize\bfseries, text=sourceblue!90!black] at (-4.4, 1.8) {Class $[a]$};
  \node[point, label={left:\scriptsize $a$}] (a) at (-4.7, 1.0) {};
  \node[point, label={right:\scriptsize $b$}] (b) at (-3.9, 1.1) {};
  \node[point, label={below:\scriptsize $c$}] (c) at (-4.3, 0.2) {};
  \draw[rel] (a) -- (b);
  \draw[rel] (b) -- (c);
  \draw[rel] (a) -- (c);

  % Class 2
  \draw[partition] (-3.1, -0.4) rectangle (-1.5, 2.1);
  \node[font=\scriptsize\bfseries, text=sourceblue!90!black] at (-2.3, 1.8) {Class $[d]$};
  \node[point, label={above:\scriptsize $d$}] (d) at (-2.3, 1.1) {};
  \node[point, label={below:\scriptsize $e$}] (e) at (-2.3, 0.2) {};
  \draw[rel] (d) -- (e);

  \node[font=\footnotesize, text=gray!80!black, align=center] at (-3.4, -0.9) {
    Every element belongs to\\
    \textbf{exactly one} class: $[a] \cap [d] = \emptyset$
  };

  % Right Panel: Dropping Transitivity
  \node[font=\bfseries\small, text=trapred!90!black] at (3.4, 2.6) {Without Transitivity: Overlapping Subsets};

  % Overlapping bubbles
  \draw[draw=quotientorange!80, fill=quotientorange!10, rounded corners=6pt, thick] (1.0, -0.4) rectangle (3.85, 2.1);
  \node[font=\tiny\bfseries, text=quotientorange!90!black] at (1.95, 1.82) {$x$'s neighbors};

  \draw[draw=trapred!70, fill=trapred!8, rounded corners=6pt, thick, opacity=0.7] (2.95, -0.4) rectangle (5.8, 2.1);
  \node[font=\tiny\bfseries, text=trapred!90!black] at (4.85, 1.82) {$z$'s neighbors};

  \node[badpoint, label={left:\scriptsize $x$}] (x) at (1.9, 0.6) {};
  \node[badpoint, label={below:\scriptsize $y$}] (y) at (3.4, 0.6) {};
  \node[badpoint, label={right:\scriptsize $z$}] (z) at (4.9, 0.6) {};

  \draw[<->, thick, quotientorange!90!black] (x) -- (y);
  \node[above=2pt, font=\tiny, text=quotientorange!90!black] at (2.2, 0.6) {$x \sim y$};

  \draw[<->, thick, trapred!90!black] (y) -- (z);
  \node[above=2pt, font=\tiny, text=trapred!90!black] at (4.6, 0.6) {$y \sim z$};

  \draw[failrel] (x) to[out=55, in=125] node[above=2pt, font=\tiny\bfseries, text=trapred] {$x \nsim z$} (z);

  \node[font=\footnotesize, text=gray!80!black, align=center] at (3.4, -0.9) {
    $y$ is shared, but $x$ and $z$ unrelated.\\
    \textbf{No partition into disjoint classes!}
  };

\end{tikzpicture}\\[0.4em]
{\small\textbf{Figure 29.3:} Equivalence partitions versus tolerance relations. On the left, transitivity ensures that elements sharing a connection form mutually disjoint equivalence classes. On the right, dropping transitivity allows relations like ``within distance $1$'' ($|x-y| \le 1$) to link $x \sim y$ and $y \sim z$ without linking $x \sim z$, producing overlapping subsets that cannot partition the set.}
\end{center}

\begin{center}
\begin{tikzpicture}[
  >=Latex,
  classbox/.style={draw=#1!80!black, fill=#1!10, rounded corners=5pt, thick},
  point/.style={circle, fill=#1!90!black, inner sep=1.8pt},
  quotpoint/.style={circle, fill=#1!90!black, inner sep=2.5pt}
]
  % Title
  \node[font=\bfseries, text=structurepurple] at (0, 3.4) {The Quotient Construction: Collapsing Classes into Single Points};

  % Original Universe S container
  \draw[draw=gray!40, fill=gray!3, rounded corners=8pt] (-5.8, 0.3) rectangle (5.8, 2.85);
  \node[font=\scriptsize\bfseries, text=gray!70!black] at (-4.7, 2.6) {Original Universe $S$};

  % Class 1: [1/2]
  \draw[classbox=sourceblue] (-5.3, 0.5) rectangle (-2.5, 2.45);
  \node[font=\scriptsize\bfseries, text=sourceblue!90!black] at (-3.9, 2.15) {Class $[1/2]$};
  \node[point=sourceblue, label={above:\tiny $2/4$}] (c1b) at (-3.9, 1.45) {};
  \node[point=sourceblue, label={below:\tiny $1/2$}] (c1a) at (-4.7, 0.95) {};
  \node[point=sourceblue, label={below:\tiny $3/6$}] (c1c) at (-3.1, 0.95) {};
  \draw[<->, thin, sourceblue!70!black] (c1a) -- (c1b);
  \draw[<->, thin, sourceblue!70!black] (c1b) -- (c1c);
  \draw[<->, thin, sourceblue!70!black] (c1a) -- (c1c);

  % Class 2: [1/3]
  \draw[classbox=targetgreen] (-1.4, 0.5) rectangle (1.4, 2.45);
  \node[font=\scriptsize\bfseries, text=targetgreen!90!black] at (0, 2.15) {Class $[1/3]$};
  \node[point=targetgreen, label={above:\tiny $3/9$}] (c2c) at (0, 1.45) {};
  \node[point=targetgreen, label={below:\tiny $1/3$}] (c2a) at (-0.8, 0.95) {};
  \node[point=targetgreen, label={below:\tiny $2/6$}] (c2b) at (0.8, 0.95) {};
  \draw[<->, thin, targetgreen!70!black] (c2a) -- (c2b);
  \draw[<->, thin, targetgreen!70!black] (c2b) -- (c2c);
  \draw[<->, thin, targetgreen!70!black] (c2a) -- (c2c);

  % Class 3: [2/3]
  \draw[classbox=quotientorange] (2.5, 0.5) rectangle (5.3, 2.45);
  \node[font=\scriptsize\bfseries, text=quotientorange!90!black] at (3.9, 2.15) {Class $[2/3]$};
  \node[point=quotientorange, label={above:\tiny $6/9$}] (c3c) at (3.9, 1.45) {};
  \node[point=quotientorange, label={below:\tiny $2/3$}] (c3a) at (3.1, 0.95) {};
  \node[point=quotientorange, label={below:\tiny $4/6$}] (c3b) at (4.7, 0.95) {};
  \draw[<->, thin, quotientorange!70!black] (c3a) -- (c3b);
  \draw[<->, thin, quotientorange!70!black] (c3b) -- (c3c);
  \draw[<->, thin, quotientorange!70!black] (c3a) -- (c3c);

  % Projection arrows
  \draw[->, very thick, sourceblue!80!black] (-3.9, 0.2) -- (-3.9, -0.6);
  \draw[->, very thick, targetgreen!80!black] (0, 0.2) -- (0, -0.6) node[midway, right=2pt, font=\scriptsize] {$\pi(x) = [x]$};
  \draw[->, very thick, quotientorange!80!black] (3.9, 0.2) -- (3.9, -0.6);

  % Quotient Set container
  \draw[draw=structurepurple!50, fill=structurepurple!5, rounded corners=8pt] (-5.8, -1.8) rectangle (5.8, -0.7);
  \node[font=\scriptsize\bfseries, text=structurepurple!90!black] at (-4.6, -0.95) {Quotient Set $S/{\sim}$};

  % Collapsed atomic points
  \node[quotpoint=sourceblue] at (-3.9, -1.3) {};
  \node[below, font=\scriptsize\bfseries, text=sourceblue!90!black] at (-3.9, -1.35) {$[1/2] = 0.5$};

  \node[quotpoint=targetgreen] at (0, -1.3) {};
  \node[below, font=\scriptsize\bfseries, text=targetgreen!90!black] at (0, -1.35) {$[1/3] \approx 0.33$};

  \node[quotpoint=quotientorange] at (3.9, -1.3) {};
  \node[below, font=\scriptsize\bfseries, text=quotientorange!90!black] at (3.9, -1.35) {$[2/3] \approx 0.67$};

\end{tikzpicture}\\[0.4em]
{\small\textbf{Figure 29.4:} The quotient set $S/{\sim}$ collapsing equivalence classes into single points. In the raw universe of fractions (top), infinitely many distinct symbols represent the same value, grouped by the equivalence relation $a/b \sim c/d \iff ad = bc$. The canonical projection map $\pi(x) = [x]$ maps each equivalence class to a single element in the quotient set $S/{\sim} = \Q$ (bottom), transforming classes into well-defined numbers.}
\end{center}

\begin{labbox}
Take the definition of an algebraic group $(G, *)$ from Chapter~21. What happens if you drop the inverse axiom? You get a \textbf{monoid}---a set with an associative operation and an identity element, but no guarantee that operations can be undone (such as natural numbers under addition, $(\N, +)$, where you cannot add a positive number to undo an addition). What happens if you drop associativity as well? You get a \textbf{magma}---an algebraic structure where elements can be combined, but grouping matters (such as integer subtraction, where $(8 - 4) - 2 = 2$ while $8 - (4 - 2) = 6$). Testing each boundary clarifies the role of each axiom in the definition of a group.
\end{labbox}

\begin{historybox}
In 1891, while waiting for a train in a Berlin railway station after a contentious mathematics conference, David Hilbert famously remarked: ``One must be able to say at all times---instead of points, lines, and planes---tables, chairs, and beer mugs.'' Hilbert was asserting that geometry is not an empirical science of physical space. A point has no intrinsic material reality. It does not matter whether a ``point'' is a speck of ink, an ordered pair of numbers, or a porcelain beer mug; what matters are purely the relations of incidence (which points lie on which lines), betweenness (which point lies between two others), and congruence that the objects satisfy with each other.

Armed with emerging symbolic logic and axiomatic set theory, Hilbert published his \emph{Grundlagen der Geometrie} in 1899, stripping geometry of sensory intuition.

This perspective led to the Frege-Hilbert correspondence between 1899 and 1903. The German logician Gottlob Frege criticized Hilbert's approach. Frege insisted on the traditional view of definitions: a definition must describe an already existing, intuitively grounded concept whose meaning is fixed before any axioms are stated. To Frege, asserting axioms about undefined terms was an ungrounded game with empty words. Hilbert countered with the modern doctrine of \textbf{implicit definition}: mathematical objects are defined entirely by the network of formal axioms they satisfy. Whether we label the primitives points and lines or tables and beer mugs, any collection of objects satisfying the axioms constitutes a legitimate geometry.

Hilbert's axiomatic method became central to twentieth-century mathematics. When mathematical objects are defined by their axiomatic properties rather than physical models, proofs become broadly applicable. An abstract algebraic theorem about groups applies simultaneously to geometric rotations, permutations of polynomial roots (Galois theory, which explains why general fifth-degree polynomials have no radical formula), and modern cryptographic keys.

Formal axiomatics also raised foundational questions. Hilbert hoped that every consistent mathematical system could prove its own consistency and completeness using purely internal symbolic mechanics (Hilbert's Program). In 1931, Kurt G\"odel showed with his Incompleteness Theorems that no consistent formal system containing basic arithmetic can prove its own consistency from within. Today, in interactive theorem provers such as Lean, Coq, and Isabelle, mathematicians face the ongoing challenge of crafting definitions that are transparent to human readers while remaining computationally sound.
\end{historybox}

\subsection{Relational definitions---and the human reader}

Reading a definition by its four parts---ambient universe, input types, defining condition, and what you get to use afterward---takes the guesswork out of both proving and applying it. And deleting one clause at a time shows why each condition is there in the first place, such as why dropping transitivity immediately ruins set partitions.

Yet axiomatic definitions carry a practical challenge: without physical pictures, proofs can become dense sequences of symbolic manipulations. When writing a proof for human readers, you must introduce variables cleanly, keep notation consistent, and manage the scope of hypotheses so the reader can follow the argument.

To see how the structure of a proof mirrors the structure of an executable program, Chapter~30 presents \textbf{Second Pass: Proofs as Programs With Human Readers}, exploring the Curry-Howard correspondence and mathematical scope management.

\begin{walkbox}
On your next walk, notice everyday definitions around you: a ``traffic stoplight,'' an ``academic degree,'' or a ``legal contract.'' Pick one and break it down into its four anatomical parts: what is the ambient universe, what is the input signature, what is the exact defining contract, and what downstream rights does it grant? What breaks down if you drop one of the requirements?
\end{walkbox}

\subsection{Exercises}

\subsubsection*{Part A: Quick Mental Checks}
\begin{enumerate}
\item For each of the following definitions, identify the \textbf{ambient universe}, the \textbf{input signature}, the \textbf{defining contract}, and the \textbf{downstream rights}:
\begin{enumerate}[label=(\alph*)]
\item A real sequence $(x_n)_{n=1}^\infty$ is \emph{bounded} if there exists a real number $M > 0$ such that $|x_n| \le M$ for all $n \in \N$.
\item A subset $U \subseteq \R$ is \emph{open} if for every $x \in U$, there exists $\eps > 0$ such that $(x - \eps, x + \eps) \subseteq U$.
\end{enumerate}

\item A student attempts to prove that $f: \R \to \R$ defined by $f(x) = 3x + 4$ is injective by writing:
\begin{quote}
``Assume $x_1 \neq x_2$. Then $3x_1 \neq 3x_2$ because $3 \neq 0$. Adding $4$ to unequal numbers gives unequal results, so $3x_1 + 4 \neq 3x_2 + 4$. Thus $f(x_1) \neq f(x_2)$.''
\end{quote}
While the student's arithmetic happens to work for this simple linear function, explain why this strategy fails for more complex functions such as $g(x) = x^3 - 3x$, and why setting $g(x_1) = g(x_2)$ provides immediate algebraic traction that inequalities cannot.

\item In the definition of an equivalence relation, what happens if we drop \textbf{reflexivity}, keeping only symmetry and transitivity? A student argues: ``If $x \sim y$, then by symmetry $y \sim x$. Then by transitivity, $(x \sim y \land y \sim x) \implies x \sim x$, so reflexivity is redundant!'' Find the subtle flaw in the student's deduction.

\item In Euclid's first proposition, an equilateral triangle is constructed on a line segment $AB$ by drawing a circle centered at $A$ of radius $AB$ and a circle centered at $B$ of radius $AB$, and picking their intersection point $C$. What unstated physical assumption was Euclid making about the two circles?
\end{enumerate}

\subsubsection*{Part B: Dissecting and Stress-Testing Definitions}
\begin{enumerate}[resume]
\item A function $f: A \to B$ is \emph{surjective} (onto) if for every $b \in B$, there exists $a \in A$ such that $f(a) = b$.
\begin{enumerate}[label=(\alph*)]
\item Dissect this definition into its four anatomical parts.
\item What downstream rights do you gain when a proof hands you the hypothesis ``$f$ is surjective''?
\item What condition must you verify if your goal is to prove ``$f$ is surjective''?
\end{enumerate}

\item Consider the relation $\le$ on the set of real numbers $\R$.
\begin{enumerate}[label=(\alph*)]
\item Verify which of the three equivalence relation axioms (reflexivity, symmetry, transitivity) hold for $\le$, and which fail.
\item Because symmetry fails, the relation is a \emph{partial order} rather than an equivalence relation. Explain geometrically why $\le$ cannot be used to partition $\R$ into disjoint clusters.
\end{enumerate}

\item Use the contrapositive definition of injectivity ($f(a_1) = f(a_2) \implies a_1 = a_2$) to prove that the rational function $f: \R \setminus \{1\} \to \R$ given by
\[
f(x) = \frac{2x + 3}{x - 1}
\]
is injective. Show every algebraic step, pointing out where downstream cancellation occurs.

\item A \emph{metric space} is a set $M$ equipped with a distance function $d: M \times M \to \R$ satisfying four axioms for all $x, y, z \in M$:
\begin{enumerate}[label=(\roman*)]
\item $d(x, y) \ge 0$ (non-negativity);
\item $d(x, y) = 0 \iff x = y$ (identity of indiscernibles);
\item $d(x, y) = d(y, x)$ (symmetry);
\item $d(x, z) \le d(x, y) + d(y, z)$ (triangle inequality).
\end{enumerate}
Stress-test this definition: what happens if we drop the direction $d(x, y) = 0 \implies x = y$ from axiom (ii), allowing distinct points to have zero distance between them? (Such a function is called a \emph{pseudometric}). Give a real-world example from daily life or physics where two distinct objects have a distance of zero under a natural measurement.
\end{enumerate}

\subsubsection*{Part C: Deeper Structural Thinking}
\begin{enumerate}[resume]
\item \textbf{Axiomatic Characterization versus Set-Theoretic Construction}: In Peano arithmetic, the natural numbers are defined through five axioms specifying the behavior of $0$ and a successor function $S(n)$. In axiomatic set theory, John von Neumann represented the natural numbers as concrete sets: $0 = \emptyset$, $1 = \{0\} = \{\emptyset\}$, $2 = \{0, 1\} = \{\emptyset, \{\emptyset\}\}$, and so on.
\begin{enumerate}[label=(\alph*)]
\item In von Neumann's representation, the statement $1 \in 2$ is a true set-theoretic fact. Does the claim ``$1$ is an element of $2$'' have any meaning in elementary arithmetic?
\item Explain how Hilbert's relational view of definitions protects working mathematicians from asking meaningless questions about the particular set-theoretic implementation of mathematical objects.
\end{enumerate}

\item \textbf{The Non-Empty Universe Clause}: In abstract algebra, a \emph{group} $(G, *)$ is defined as a set $G$ with a binary operation $*$ satisfying associativity, containing an identity element $e \in G$, and providing an inverse $a^{-1} \in G$ for every element $a \in G$. Almost every textbook includes the upfront ambient clause: ``$G$ is a non-empty set.''
\begin{enumerate}[label=(\alph*)]
\item What happens to the identity axiom if $G = \emptyset$? Can an empty set contain an identity element?
\item If an author omitted the non-empty clause and relaxed the identity axiom so that the empty set counted as a group, what caveat would have to be attached to almost every theorem in group theory (such as Lagrange's theorem on the order of subgroups)?
\end{enumerate}
\end{enumerate}

\newpage
